\section{Stability Filtering}

Stability filtering ensures that semantic states evolve smoothly over time,
preventing jitter, noise amplification, and abrupt geometric transitions.
In the MEM8-Geometric-Architecture, stability is implemented using a Kalman
filter, which provides optimal state estimation under uncertainty.

Let $\mathbf{x}_t \in \mathbb{R}^d$ denote the fused semantic state at time $t$,
and let $\mathbf{u}_t$ represent optional control inputs. The stability filter
operates in two phases: prediction and update.

\subsection{Prediction Step}

The predicted semantic state is given by:
\[
\mathbf{x}_{t|t-1} = A \mathbf{x}_{t-1} + B \mathbf{u}_t
\]

Here,
\[
A \in \mathbb{R}^{d \times d}, \quad B \in \mathbb{R}^{d \times m}
\]
are learned geometric transition matrices. The predicted covariance is:
\[
P_{t|t-1} = A P_{t-1} A^\top + Q
\]
where $Q$ is the process noise covariance.

\subsection{Update Step}

Given an observation $\mathbf{y}_t$, the updated semantic state is:
\[
\mathbf{x}_{t|t} =
\mathbf{x}_{t|t-1} +
K_t (\mathbf{y}_t - \mathbf{x}_{t|t-1})
\]

The Kalman gain is:
\[
K_t =
P_{t|t-1} H^\top
\left(
H P_{t|t-1} H^\top + R
\right)^{-1}
\]

The updated covariance is:
\[
P_{t|t} = (I - K_t H) P_{t|t-1}
\]

\subsection{Geometric Interpretation}

Stability filtering can be viewed as a curvature-aware smoothing operation on
the semantic manifold. The prediction step advances the state along a learned
geometric trajectory, while the update step corrects deviations using observed
data.

The Kalman gain determines how strongly the system trusts the observation
versus the predicted trajectory. High gain corresponds to rapid correction,
while low gain corresponds to smooth continuation along the manifold.

This geometric smoothing is essential for preventing:

\begin{itemize}
\item drift amplification,
\item curvature spikes,
\item collapse-risk escalation,
\item abrupt semantic jumps,
\item instability in downstream routing.
\end{itemize}

\subsection{Role in the Cognitive Loop}

Stability filtering operates after salience detection and before semantic
routing. High-salience states may trigger stronger updates, while low-salience
states rely more heavily on predicted trajectories.

This ensures that the cognitive loop remains both responsive and stable,
balancing novelty with continuity.
